% Generated draft: requires mathematical and source review.
\documentclass[11pt]{article}
\usepackage[margin=25mm]{geometry}
\usepackage{fontspec}
\setmainfont{Latin Modern Roman}
\usepackage{amsmath,amssymb,mathtools,unicode-math}
\setmathfont{Latin Modern Math}
\usepackage{enumitem,needspace,xurl,hyperref}
\setlength{\parindent}{0pt}
\setlength{\parskip}{5pt}
\setlength{\emergencystretch}{3em}
\newcommand{\N}{\mathbb{N}}
\newcommand{\Z}{\mathbb{Z}}
\newcommand{\Q}{\mathbb{Q}}
\newcommand{\R}{\mathbb{R}}
\newcommand{\C}{\mathbb{C}}
\newcommand{\F}{\mathbb{F}}
\newcommand{\A}{\mathbb{A}}
\newcommand{\PP}{\mathbb{P}}
\newcommand{\E}{\mathbb{E}}
\newcommand{\eps}{\varepsilon}
\newcommand{\dd}{\,\mathrm{d}}
\begin{document}
\section*{UnsolvedMath ID: 800515}
\textbf{Source problem code:} catalogue entry
\subsection{Definitions and mathematical statement}
Let \(n\geq 1\) and \(A\in\C^{n\times n}\). For a nonzero starting vector \(b\in\C^n\) and an integer \(1\leq m\leq n\), define the Krylov space
\[
\mathcal K_m(A,b)=\operatorname{span}\{b,Ab,\ldots,A^{m-1}b\}.
\]
Write \(d=\dim\mathcal K_m(A,b)\), and let \(Q\in\C^{n\times d}\) satisfy \(Q^*Q=I_d\) and \(\operatorname{range}(Q)=\mathcal K_m(A,b)\). Set \(H=Q^*AQ\). The source specifies a random starting vector but does not specify its probability law, the Krylov order \(m\), or how to handle a zero starting vector. The actual dimension \(d\) allows for linear dependence among the Krylov vectors.

Under the usual eigenvector condition-number interpretation, define
\[
\kappa_V(H)=\inf\bigl\{\|V\|_2\|V^{-1}\|_2: V\in\C^{d\times d}\text{ is invertible and }V^{-1}HV\text{ is diagonal}\bigr\},
\]
with \(\kappa_V(H)=+\infty\) if \(H\) is not diagonalizable. Here \(\|\cdot\|_2\) is the operator norm induced by the Euclidean vector norm. This convention for \(\kappa_V\) is a standard interpretation of notation that the supplied record does not define explicitly.

For \(\varepsilon>0\), use the spectral-norm pseudospectrum
\[
\Lambda_\varepsilon(H)=\{z\in\C:\sigma_{\min}(zI_d-H)\leq\varepsilon\},
\]
where \(\sigma_{\min}\) denotes the smallest singular value. Let \(\operatorname{area}\) denote two-dimensional Lebesgue measure on \(\C\cong\R^2\). The pseudospectral norm convention is also not specified in the source.

The source asks whether, for arbitrary \(A\), random Krylov compression yields a polynomially bounded eigenvector condition number with high probability, or whether it yields an expected pseudospectral area bounded by a polynomial in \(n\) times \(\varepsilon^\beta\), for \(\beta\) close to \(2\). Schematically, these proposed bounds have the forms
\[
\Pr_b\!\left\{\kappa_V(Q^*AQ)\leq p(n)\right\}\geq 1-\delta_n
\qquad\text{or}\qquad
\E_b\!\left[\operatorname{area}\Lambda_\varepsilon(Q^*AQ)\right]\leq p(n)\varepsilon^\beta,
\]
where \(p\) is a polynomial. These displays are formulation templates: the source does not specify the law of \(b\), the permitted Krylov orders, the failure probability \(\delta_n\), the range of \(\varepsilon\), the precise meaning of \(\beta\) being close to \(2\), or the permitted dependence of constants and exponents on these parameters or on \(A\). No normalization of \(A\) is supplied. Consequently, the record does not determine a fully quantified conjecture. The source additionally reports that the problem is open even when \(A\) is a circulant shift matrix; this is a supplied status qualification, not an independently verified literature assessment.
\subsection{Short English statement}
For an arbitrary complex matrix, form its compression to a Krylov space generated by a random vector. The source asks whether the compressed matrix has polynomially controlled eigenvector conditioning with high probability, or an expected \(\varepsilon\)-pseudospectral area at most \(\mathrm{poly}(n)\varepsilon^\beta\) with \(\beta\) close to \(2\). It reports openness even for a circulant shift matrix. The random-vector law, Krylov order, quantitative probability threshold, and other parameter qualifications remain unspecified.
\subsection{Sources}
\begin{itemize}
\item Supplied local UnsolvedMath registry record, permanent ID \(800515\), exact source problem code \texttt{catalogue entry}, title ``Linear Systems and Eigenvalue Problems --- Problem 24.'' The registry attributes the problem to Amsel, Baumann et al., proposed year 2026, and labels it open. These are supplied attribution and status data.
\item Amsel, Baumann et al., \emph{Linear Systems and Eigenvalue Problems: Open Questions from a Simons Workshop} (2026), Problem 24. Supplied source URL: \url{https://arxiv.org/abs/2602.05394}. The record reports access on 2026-10-16 and extraction method \texttt{arxiv\_latex}. These provenance details are supplied and have not been independently verified here.
\item Supplied upstream revision: \texttt{d371ee917ac9647fbb06cd6fc529deec6f3dda82}. The supplied statement includes the qualification that the problem is open even for a circulant shift matrix; no separate supporting citation for that qualification is supplied.
\end{itemize}
\end{document}
